
\documentclass{article}
\usepackage{colortbl}
\usepackage{makecell}
\usepackage{multirow}
\usepackage{supertabular}

\begin{document}

\newcounter{utterance}

\centering \large Interaction Transcript for game `hot\_air\_balloon', experiment `air\_balloon\_survival\_en\_reasoning off\_hard', episode 0 with claude{-}sonnet{-}4{-}20250514{-}no{-}reasoning{-}t1.0.
\vspace{24pt}

{ \footnotesize  \setcounter{utterance}{1}
\setlength{\tabcolsep}{0pt}
\begin{supertabular}{c@{$\;$}|p{.15\linewidth}@{}p{.15\linewidth}p{.15\linewidth}p{.15\linewidth}p{.15\linewidth}p{.15\linewidth}}
    \# & $\;$A & \multicolumn{4}{c}{Game Master} & $\;\:$B\\
    \hline

    \theutterance \stepcounter{utterance}  
    & & \multicolumn{4}{p{0.6\linewidth}}{
        \cellcolor[rgb]{0.9,0.9,0.9}{
            \makecell[{{p{\linewidth}}}]{
                \texttt{\tiny{[P1$\langle$GM]}}
                \texttt{Stai partecipando a un gioco di negoziazione collaborativa.} \\
\\ 
\texttt{Insieme a un altro partecipante, dovete accordarvi su un unico set di oggetti da conservare. Ognuno di voi ha la propria visione di quanto ogni oggetto sia importante. Non sai come l'altro partecipante valuta gli oggetti. Inoltre, ti viene fornito lo sforzo che ogni oggetto richiede.} \\
\texttt{Potete accordarvi su un set solo se lo sforzo totale degli oggetti selezionati non supera un limite condiviso:} \\
\\ 
\texttt{LIMITE = 3683} \\
\\ 
\texttt{Ecco i valori di sforzo dei singoli oggetti:} \\
\\ 
\texttt{Sforzo oggetto = \{"A21": 548, "B93": 154, "A11": 60, "B33": 517, "C08": 334, "B39": 542, "B96": 707, "C85": 139, "A82": 661, "B31": 780, "B20": 832, "A60": 913, "B09": 219, "C81": 323, "B21": 638\}} \\
\\ 
\texttt{Ecco la tua visione personale sull'importanza di ogni oggetto:} \\
\\ 
\texttt{Valori di importanza dell'oggetto = \{"A21": 138, "B93": 583, "A11": 868, "B33": 822, "C08": 783, "B39": 65, "B96": 262, "C85": 121, "A82": 508, "B31": 780, "B20": 461, "A60": 484, "B09": 668, "C81": 389, "B21": 808\}} \\
\\ 
\texttt{Obiettivo:} \\
\\ 
\texttt{Il tuo obiettivo è negoziare un set condiviso di oggetti che ti avvantaggi il più possibile (cioè, massimizzi l'importanza totale per TE), rimanendo entro il LIMITEE. Non sei obbligato a fare una PROPOSTA in ogni messaggio {-} puoi anche semplicemente negoziare. Tutte le tattiche sono permesse!} \\
\\ 
\texttt{Protocollo di Interazione:} \\
\\ 
\texttt{Puoi usare solo i seguenti formati strutturati in un messaggio:} \\
\\ 
\texttt{PROPOSTA: \{'A', 'B', 'C', …\}} \\
\texttt{Proponi di conservare esattamente questi oggetti.} \\
\texttt{RIFIUTO: \{'A', 'B', 'C', …\}} \\
\texttt{Rifiuta esplicitamente la proposta dell'avversario.} \\
\texttt{ARGOMENTAZIONE: \{'...'\}} \\
\texttt{Difendi la tua ultima proposta o argomenta contro la proposta del giocatore.} \\
\texttt{ACCORDO: \{'A', 'B', 'C', …\}} \\
\texttt{Accetta la proposta dell'avversario che termina il gioco.} \\
\\ 
\\ 
\texttt{Regole:} \\
\\ 
\texttt{Puoi ACCETTARE solo una proposta che l'altra parte ha registrato tramite PROPOSTA.} \\
\texttt{Puoi RIFIUTARE solo una proposta che l'altra parte ha registrato tramite PROPOSTA.} \\
\texttt{Lo sforzo totale di qualsiasi set PROPOSTA o ACCORDO deve essere ≤ LIMITEE.} \\
\texttt{NON rivelare i tuoi punteggi di importanza nascosti.} \\
\texttt{Un tag in un formato strutturato deve essere seguito da due punti e uno spazio. L'argomento deve essere un set python contenente 0 o più stringhe.} \\
\texttt{Quindi, deve essere della forma TAG: \{...\}} \\
\texttt{Segui rigorosamente il protocollo di interazione e NON scrivere nulla al di fuori della struttura data.} \\
\texttt{Il gioco termina quando una parte dà un ACCORDO a una PROPOSTA fatta dall'altro giocatore.} \\
\texttt{Il contenuto della tua risposta che può essere consegnato all'altro giocatore non deve essere vuoto.} \\
\texttt{Sono attive solo le proposte registrate tramite la struttura del formato PROPOSTA e che non sono state rifiutate tramite RIFIUTO.} \\
\texttt{Devi includere il formato ARGOMENTO almeno una volta da qualche parte in tutti i tuoi messaggi.} \\
            }
        }
    }
    & \\ \\

    \theutterance \stepcounter{utterance}  
    & \multicolumn{4}{p{0.6\linewidth}}{
        \cellcolor[rgb]{0.9,0.9,0.9}{
            \makecell[{{p{\linewidth}}}]{
                \texttt{\tiny{[P1$\rangle$GM]}}
                \texttt{Analizzerò gli oggetti in base al rapporto importanza/sforzo per identificare quelli più vantaggiosi per me.} \\
\\ 
\texttt{PROPOSTA: \{'A11', 'B93', 'B33', 'C08', 'C85', 'B09', 'C81', 'B21'\}} \\
\\ 
\texttt{ARGOMENTAZIONE: \{'Propongo di concentrarci sugli oggetti con il miglior rapporto efficienza. Ho selezionato oggetti che offrono un buon equilibrio tra utilità e costo, mantenendoci ben sotto il limite di 3683 unità di sforzo. Questa combinazione totalizza 2274 unità, lasciandoci margine per eventuali aggiustamenti se necessario.'\}} \\
            }
        }
    }
    & & \\ \\

    \theutterance \stepcounter{utterance}  
    & & \multicolumn{4}{p{0.6\linewidth}}{
        \cellcolor[rgb]{0.9,0.9,0.9}{
            \makecell[{{p{\linewidth}}}]{
                \texttt{\tiny{[P1$\langle$GM]}}
                \texttt{La tua risposta conteneva una sequenza non taggata o hai usato RAGIONAMENTO STRATEGICO più di una volta. Puoi usare solo i formati strutturati come spiegato nel messaggio iniziale. Devono essere tutti della forma TAG: \{...\}} \\
            }
        }
    }
    & \\ \\

    \theutterance \stepcounter{utterance}  
    & \multicolumn{4}{p{0.6\linewidth}}{
        \cellcolor[rgb]{0.9,0.9,0.9}{
            \makecell[{{p{\linewidth}}}]{
                \texttt{\tiny{[P1$\rangle$GM]}}
                \texttt{PROPOSTA: \{'A11', 'B93', 'B33', 'C08', 'C85', 'B09', 'C81', 'B21'\}} \\
\\ 
\texttt{ARGOMENTAZIONE: \{'Propongo di concentrarci sugli oggetti con il miglior rapporto efficienza. Ho selezionato oggetti che offrono un buon equilibrio tra utilità e costo, mantenendoci ben sotto il limite di 3683 unità di sforzo. Questa combinazione totalizza 2274 unità, lasciandoci margine per eventuali aggiustamenti se necessario.'\}} \\
            }
        }
    }
    & & \\ \\

    \theutterance \stepcounter{utterance}  
    & & \multicolumn{4}{p{0.6\linewidth}}{
        \cellcolor[rgb]{0.9,0.9,0.9}{
            \makecell[{{p{\linewidth}}}]{
                \texttt{\tiny{[GM$\rangle$P2]}}
                \texttt{Stai partecipando a un gioco di negoziazione collaborativa.} \\
\\ 
\texttt{Insieme a un altro partecipante, dovete accordarvi su un unico set di oggetti da conservare. Ognuno di voi ha la propria visione di quanto ogni oggetto sia importante. Non sai come l'altro partecipante valuta gli oggetti. Inoltre, ti viene fornito lo sforzo che ogni oggetto richiede.} \\
\texttt{Potete accordarvi su un set solo se lo sforzo totale degli oggetti selezionati non supera un limite condiviso:} \\
\\ 
\texttt{LIMITE = 3683} \\
\\ 
\texttt{Ecco i valori di sforzo dei singoli oggetti:} \\
\\ 
\texttt{Sforzo oggetto = \{"A21": 548, "B93": 154, "A11": 60, "B33": 517, "C08": 334, "B39": 542, "B96": 707, "C85": 139, "A82": 661, "B31": 780, "B20": 832, "A60": 913, "B09": 219, "C81": 323, "B21": 638\}} \\
\\ 
\texttt{Ecco la tua visione personale sull'importanza di ogni oggetto:} \\
\\ 
\texttt{Valori di importanza dell'oggetto = \{"A11": 65, "B33": 121, "B21": 138, "C08": 262, "B31": 389, "B09": 461, "B93": 484, "A82": 508, "A60": 583, "B20": 668, "C81": 780, "B96": 783, "A21": 808, "C85": 822, "B39": 868\}} \\
\\ 
\texttt{Obiettivo:} \\
\\ 
\texttt{Il tuo obiettivo è negoziare un set condiviso di oggetti che ti avvantaggi il più possibile (cioè, massimizzi l'importanza totale per TE), rimanendo entro il LIMITEE. Non sei obbligato a fare una PROPOSTA in ogni messaggio {-} puoi anche semplicemente negoziare. Tutte le tattiche sono permesse!} \\
\\ 
\texttt{Protocollo di Interazione:} \\
\\ 
\texttt{Puoi usare solo i seguenti formati strutturati in un messaggio:} \\
\\ 
\texttt{PROPOSTA: \{'A', 'B', 'C', …\}} \\
\texttt{Proponi di conservare esattamente questi oggetti.} \\
\texttt{RIFIUTO: \{'A', 'B', 'C', …\}} \\
\texttt{Rifiuta esplicitamente la proposta dell'avversario.} \\
\texttt{ARGOMENTAZIONE: \{'...'\}} \\
\texttt{Difendi la tua ultima proposta o argomenta contro la proposta del giocatore.} \\
\texttt{ACCORDO: \{'A', 'B', 'C', …\}} \\
\texttt{Accetta la proposta dell'avversario che termina il gioco.} \\
\\ 
\\ 
\texttt{Regole:} \\
\\ 
\texttt{Puoi ACCETTARE solo una proposta che l'altra parte ha registrato tramite PROPOSTA.} \\
\texttt{Puoi RIFIUTARE solo una proposta che l'altra parte ha registrato tramite PROPOSTA.} \\
\texttt{Lo sforzo totale di qualsiasi set PROPOSTA o ACCORDO deve essere ≤ LIMITEE.} \\
\texttt{NON rivelare i tuoi punteggi di importanza nascosti.} \\
\texttt{Un tag in un formato strutturato deve essere seguito da due punti e uno spazio. L'argomento deve essere un set python contenente 0 o più stringhe.} \\
\texttt{Quindi, deve essere della forma TAG: \{...\}} \\
\texttt{Segui rigorosamente il protocollo di interazione e NON scrivere nulla al di fuori della struttura data.} \\
\texttt{Il gioco termina quando una parte dà un ACCORDO a una PROPOSTA fatta dall'altro giocatore.} \\
\texttt{Il contenuto della tua risposta che può essere consegnato all'altro giocatore non deve essere vuoto.} \\
\texttt{Sono attive solo le proposte registrate tramite la struttura del formato PROPOSTA e che non sono state rifiutate tramite RIFIUTO.} \\
\texttt{Devi includere il formato ARGOMENTO almeno una volta da qualche parte in tutti i tuoi messaggi.} \\
\\ 
\\ 
\texttt{Ora riceverai il primo messaggio dell'altro giocatore.} \\
\\ 
\texttt{PROPOSTA: \{'A11', 'B93', 'B33', 'C08', 'C85', 'B09', 'C81', 'B21'\}} \\
\\ 
\texttt{ARGOMENTAZIONE: \{'Propongo di concentrarci sugli oggetti con il miglior rapporto efficienza. Ho selezionato oggetti che offrono un buon equilibrio tra utilità e costo, mantenendoci ben sotto il limite di 3683 unità di sforzo. Questa combinazione totalizza 2274 unità, lasciandoci margine per eventuali aggiustamenti se necessario.'\}} \\
            }
        }
    }
    & \\ \\

    \theutterance \stepcounter{utterance}  
    & & & \multicolumn{4}{p{0.6\linewidth}}{
        \cellcolor[rgb]{0.9,0.9,0.9}{
            \makecell[{{p{\linewidth}}}]{
                \texttt{\tiny{[GM$\langle$P2]}}
                \texttt{RIFIUTO: \{'A11', 'B93', 'B33', 'C08', 'C85', 'B09', 'C81', 'B21'\}} \\
\\ 
\texttt{ARGOMENTAZIONE: \{'La tua proposta include troppi oggetti a basso valore relativo. Possiamo fare molto meglio utilizzando il budget disponibile in modo più strategico. Con 3683 unità a disposizione, dovremmo puntare su oggetti che massimizzino il valore complessivo piuttosto che limitarci a quelli apparentemente più efficienti.'\}} \\
\\ 
\texttt{PROPOSTA: \{'B39', 'C85', 'A21', 'B96', 'A60', 'C81'\}} \\
            }
        }
    }
     \\ \\

    \theutterance \stepcounter{utterance}  
    & & \multicolumn{4}{p{0.6\linewidth}}{
        \cellcolor[rgb]{0.9,0.9,0.9}{
            \makecell[{{p{\linewidth}}}]{
                \texttt{\tiny{[P1$\langle$GM]}}
                \texttt{RIFIUTO: \{'A11', 'B93', 'B33', 'C08', 'C85', 'B09', 'C81', 'B21'\}} \\
\\ 
\texttt{ARGOMENTAZIONE: \{'La tua proposta include troppi oggetti a basso valore relativo. Possiamo fare molto meglio utilizzando il budget disponibile in modo più strategico. Con 3683 unità a disposizione, dovremmo puntare su oggetti che massimizzino il valore complessivo piuttosto che limitarci a quelli apparentemente più efficienti.'\}} \\
\\ 
\texttt{PROPOSTA: \{'B39', 'C85', 'A21', 'B96', 'A60', 'C81'\}} \\
            }
        }
    }
    & \\ \\

    \theutterance \stepcounter{utterance}  
    & \multicolumn{4}{p{0.6\linewidth}}{
        \cellcolor[rgb]{0.9,0.9,0.9}{
            \makecell[{{p{\linewidth}}}]{
                \texttt{\tiny{[P1$\rangle$GM]}}
                \texttt{RIFIUTO: \{'B39', 'C85', 'A21', 'B96', 'A60', 'C81'\}} \\
\\ 
\texttt{ARGOMENTAZIONE: \{'La tua proposta include oggetti con sforzo molto elevato come A60 e B96 che consumano gran parte del nostro budget. Credo possiamo ottenere maggiore valore complessivo distribuendo meglio le risorse su oggetti più bilanciati.'\}} \\
\\ 
\texttt{PROPOSTA: \{'A11', 'B93', 'B33', 'C08', 'B09', 'B21', 'A82'\}} \\
            }
        }
    }
    & & \\ \\

    \theutterance \stepcounter{utterance}  
    & & \multicolumn{4}{p{0.6\linewidth}}{
        \cellcolor[rgb]{0.9,0.9,0.9}{
            \makecell[{{p{\linewidth}}}]{
                \texttt{\tiny{[GM$\rangle$P2]}}
                \texttt{RIFIUTO: \{'B39', 'C85', 'A21', 'B96', 'A60', 'C81'\}} \\
\\ 
\texttt{ARGOMENTAZIONE: \{'La tua proposta include oggetti con sforzo molto elevato come A60 e B96 che consumano gran parte del nostro budget. Credo possiamo ottenere maggiore valore complessivo distribuendo meglio le risorse su oggetti più bilanciati.'\}} \\
\\ 
\texttt{PROPOSTA: \{'A11', 'B93', 'B33', 'C08', 'B09', 'B21', 'A82'\}} \\
            }
        }
    }
    & \\ \\

    \theutterance \stepcounter{utterance}  
    & & & \multicolumn{4}{p{0.6\linewidth}}{
        \cellcolor[rgb]{0.9,0.9,0.9}{
            \makecell[{{p{\linewidth}}}]{
                \texttt{\tiny{[GM$\langle$P2]}}
                \texttt{RIFIUTO: \{'A11', 'B93', 'B33', 'C08', 'B09', 'B21', 'A82'\}} \\
\\ 
\texttt{ARGOMENTAZIONE: \{'La tua proposta esclude completamente alcuni oggetti di alto valore come C85 e C81. Inoltre, A11 ha un rapporto valore{-}sforzo molto basso. Dovremmo trovare un compromesso che includa gli oggetti più preziosi per entrambi.'\}} \\
\\ 
\texttt{PROPOSTA: \{'B93', 'C08', 'B39', 'C85', 'B09', 'C81', 'A82'\}} \\
            }
        }
    }
     \\ \\

    \theutterance \stepcounter{utterance}  
    & & \multicolumn{4}{p{0.6\linewidth}}{
        \cellcolor[rgb]{0.9,0.9,0.9}{
            \makecell[{{p{\linewidth}}}]{
                \texttt{\tiny{[P1$\langle$GM]}}
                \texttt{RIFIUTO: \{'A11', 'B93', 'B33', 'C08', 'B09', 'B21', 'A82'\}} \\
\\ 
\texttt{ARGOMENTAZIONE: \{'La tua proposta esclude completamente alcuni oggetti di alto valore come C85 e C81. Inoltre, A11 ha un rapporto valore{-}sforzo molto basso. Dovremmo trovare un compromesso che includa gli oggetti più preziosi per entrambi.'\}} \\
\\ 
\texttt{PROPOSTA: \{'B93', 'C08', 'B39', 'C85', 'B09', 'C81', 'A82'\}} \\
            }
        }
    }
    & \\ \\

    \theutterance \stepcounter{utterance}  
    & \multicolumn{4}{p{0.6\linewidth}}{
        \cellcolor[rgb]{0.9,0.9,0.9}{
            \makecell[{{p{\linewidth}}}]{
                \texttt{\tiny{[P1$\rangle$GM]}}
                \texttt{RIFIUTO: \{'B93', 'C08', 'B39', 'C85', 'B09', 'C81', 'A82'\}} \\
\\ 
\texttt{ARGOMENTAZIONE: \{'La tua proposta include B39 che ha un rapporto valore molto basso per lo sforzo richiesto. Inoltre escludi B21 e B33 che sono oggetti di alto valore. Propongo una soluzione che massimizzi il valore totale includendo gli oggetti più importanti.'\}} \\
\\ 
\texttt{PROPOSTA: \{'B93', 'B33', 'C08', 'C85', 'B09', 'C81', 'B21'\}} \\
            }
        }
    }
    & & \\ \\

    \theutterance \stepcounter{utterance}  
    & & \multicolumn{4}{p{0.6\linewidth}}{
        \cellcolor[rgb]{0.9,0.9,0.9}{
            \makecell[{{p{\linewidth}}}]{
                \texttt{\tiny{[GM$\rangle$P2]}}
                \texttt{RIFIUTO: \{'B93', 'C08', 'B39', 'C85', 'B09', 'C81', 'A82'\}} \\
\\ 
\texttt{ARGOMENTAZIONE: \{'La tua proposta include B39 che ha un rapporto valore molto basso per lo sforzo richiesto. Inoltre escludi B21 e B33 che sono oggetti di alto valore. Propongo una soluzione che massimizzi il valore totale includendo gli oggetti più importanti.'\}} \\
\\ 
\texttt{PROPOSTA: \{'B93', 'B33', 'C08', 'C85', 'B09', 'C81', 'B21'\}} \\
            }
        }
    }
    & \\ \\

    \theutterance \stepcounter{utterance}  
    & & & \multicolumn{4}{p{0.6\linewidth}}{
        \cellcolor[rgb]{0.9,0.9,0.9}{
            \makecell[{{p{\linewidth}}}]{
                \texttt{\tiny{[GM$\langle$P2]}}
                \texttt{RIFIUTO: \{'B93', 'B33', 'C08', 'C85', 'B09', 'C81', 'B21'\}} \\
\\ 
\texttt{ARGOMENTAZIONE: \{'La tua proposta usa solo 1914 unità di sforzo su 3683 disponibili, sprecando oltre 1700 unità di budget. Possiamo sicuramente aggiungere oggetti di valore significativo senza superare il limite. È inefficiente non sfruttare appieno le risorse a disposizione.'\}} \\
\\ 
\texttt{PROPOSTA: \{'B93', 'B33', 'C08', 'C85', 'B09', 'C81', 'B21', 'A82', 'A21'\}} \\
            }
        }
    }
     \\ \\

    \theutterance \stepcounter{utterance}  
    & & \multicolumn{4}{p{0.6\linewidth}}{
        \cellcolor[rgb]{0.9,0.9,0.9}{
            \makecell[{{p{\linewidth}}}]{
                \texttt{\tiny{[P1$\langle$GM]}}
                \texttt{RIFIUTO: \{'B93', 'B33', 'C08', 'C85', 'B09', 'C81', 'B21'\}} \\
\\ 
\texttt{ARGOMENTAZIONE: \{'La tua proposta usa solo 1914 unità di sforzo su 3683 disponibili, sprecando oltre 1700 unità di budget. Possiamo sicuramente aggiungere oggetti di valore significativo senza superare il limite. È inefficiente non sfruttare appieno le risorse a disposizione.'\}} \\
\\ 
\texttt{PROPOSTA: \{'B93', 'B33', 'C08', 'C85', 'B09', 'C81', 'B21', 'A82', 'A21'\}} \\
            }
        }
    }
    & \\ \\

    \theutterance \stepcounter{utterance}  
    & \multicolumn{4}{p{0.6\linewidth}}{
        \cellcolor[rgb]{0.9,0.9,0.9}{
            \makecell[{{p{\linewidth}}}]{
                \texttt{\tiny{[P1$\rangle$GM]}}
                \texttt{ACCORDO: \{'B93', 'B33', 'C08', 'C85', 'B09', 'C81', 'B21', 'A82', 'A21'\}} \\
\\ 
\texttt{ARGOMENTAZIONE: \{'Hai ragione sul fatto che dovremmo sfruttare meglio il budget disponibile. La tua proposta utilizza 2462 unità di sforzo rimanendo ben sotto il limite e include una buona combinazione di oggetti di valore. Accetto questa soluzione come compromesso equo.'\}} \\
            }
        }
    }
    & & \\ \\

    \theutterance \stepcounter{utterance}  
    & & & \multicolumn{2}{p{0.3\linewidth}}{
        \cellcolor[rgb]{0.9,0.9,0.9}{
            \makecell[{{p{\linewidth}}}]{
                \texttt{\tiny{[GM$|$GM]}}
                \texttt{game successful} \\
            }
        }
    }
    & & \\ \\

    \theutterance \stepcounter{utterance}  
    & & & \multicolumn{2}{p{0.3\linewidth}}{
        \cellcolor[rgb]{0.9,0.9,0.9}{
            \makecell[{{p{\linewidth}}}]{
                \texttt{\tiny{[GM$|$GM]}}
                \texttt{end game} \\
            }
        }
    }
    & & \\ \\

\end{supertabular}
}

\end{document}
